\documentclass[../../main.tex]{subfiles}

\begin{document}

\chapter{Wave Function Collapse}
Wave Function Collapse is a procedural generation algorithm that assigns tiles to cells in a grid while satisfying local compatibility rules. Since its introduction, several extensions have been proposed to address its limitations and support different applications. In this chapter, we explain how the standard WFC algorithm works and review its main variants and applications, focusing on the techniques relevant to this thesis.

\subfile{wfc-algorithm.tex}
\subfile{wfc-variants.tex}

\section{Summary}

In this chapter, we reviewed the Wave Function Collapse algorithm, its main extensions, and its applications in different domains.

\paragraph{WFC Extensions.}
WFC assigns tiles to grid cells while enforcing local compatibility rules~\citeonline{Gumin2016wfc}. Its behavior depends on the strategies used to select cells and tiles, as well as on how contradictions are handled~\citebibliography{Karth2022a}. Several extensions introduce global constraints and dynamic tile weights~\citebibliography{Sandh2019a}, interactive control~\citebibliography{Lang2022a}, graph-based structures~\citebibliography{Kim2019a, Kim2025a}, nested grids for scalability~\citebibliography{Nie2023a}, and tiles of different sizes~\citebibliography{Piep2025a}. WFC has also been applied to procedural content generation, particularly for game environments and levels~\citeonline{ParkerProcSkater, StalbergBadNorth, StalbergTownscaper, KleinebergInfiniteCity} and to poetry generation~\citeonline{OlearyPoetry}.

\paragraph{Influence on this Thesis.}
These extensions provide several ideas relevant to our approach. We adopt dynamic tile weight adjustment~\citebibliography{Sandh2019a}, backtracking strategies~\citebibliography{Karth2022a}, and non-uniform tiles~\citebibliography{Piep2025a}. However, applying WFC to software visualization requires additional mechanisms because local compatibility alone does not ensure that software relationships connect the correct entities. We therefore define tiles and compatibility rules based on software relationships, combine WFC with Multidimensional Scaling for entity placement, and introduce mechanisms to guide paths towards their destinations and handle contradictions. These adaptations let us explore WFC for drawing software relationships while preserving connectivity and reducing path crossings.

\end{document}